%% kapa: exam=1 round=28 year=2017 period=1 form=A subject=economics subject_ko=경제학원론 seq=23 number=63
%% source: Q-Net 감정평가사 기출문제 (cst00309) artlSeq=5199603 file=extracted/1008/5199603/2017년 제28회 감정평가사 1차시험 문제지/1교시 A형.pdf
%% conversion: kordoc 4.24.51 -> tools/gen1.py (auto); review=proofread-sonnet; flags=none
\begin{question}{63}
고정환율제인 먼델-플레밍 모형에서 해외이자율이 상승할 경우, 자국에 나타나는 경제변화에 관한 설명으로 옳은 것은? (단, 자국은 자본이동이 완전히 자유로운 소규모 개방경제국이다.)
\begin{choices}
  \item 환율은 불변이고, 생산량은 감소한다.
  \item 환율은 불변이고, 무역수지는 증가한다.
  \item 환율은 불변이고, 국내투자수요가 증가한다.
  \item 환율에 대한 하락압력으로 통화량이 증가한다.
  \item 국내이자율이 하락함에 따라 국내투자수요가 증가한다.
\end{choices}
\end{question}
\answer{1}
